\documentclass{article}
\usepackage{graphicx} % Required for inserting images
\usepackage{amsmath}
\usepackage{}

\title{Inecuacions}
\author{Lluis Gibert i Liam Korsak}
\date{23 de març de 2026}

\begin{document}

\maketitle

\section*{Exercici 1}



\[
\text{Resolver: } \left|x-\left|x-2\right|\right|<3
\]

Estudiam per casos:

\[
\textbf{Cas 1: } x\geq 2
\]

Llavors:

\[
|x-2|=x-2
\]

y la inecuació queda:

\[
|x-(x-2)|<3
\]

\[
|2|<3
\]

\[
2<3
\]

Això és cert, llavors:

\[
x\geq 2
\]

són solucions.

\vspace{0.5cm}

\[
\textbf{Cas 2: } x<2
\]

Llavors:

\[
|x-2|=2-x
\]

y la inecuación queda:

\[
|x-(2-x)|<3
\]

\[
|2x-2|<3
\]

\[
2|x-1|<3
\]

\[
|x-1|<\frac{3}{2}
\]

Això equival a:

\[
-\frac{3}{2}<x-1<\frac{3}{2}
\]

Sumant \(1\) en els tres membres:

\[
-\frac{1}{2}<x<\frac{5}{2}
\]

Pero en aquest cas estam suposant \(x<2\), així que:

\[
-\frac{1}{2}<x<2
\]

\vspace{0.5cm}

Unint els dos casos:

\[
\left(-\frac{1}{2},2\right)\cup [2,+\infty)
\]

\[
\boxed{\left(-\frac{1}{2},+\infty\right)}
\]

\section*{Exercici 2}
\[
x^2-ax\le 2
\]

Passem tot a un costat:
\[
x^2-ax-2\le 0
\]

Resolem l'equació associada:
\[
x^2-ax-2=0
\]

\[
x=\frac{a\pm\sqrt{a^2+8}}{2}
\]

Com que el coeficient de \(x^2\) és positiu, la paràbola obre cap amunt, i per tant
la inecuació es compleix entre les dues arrels.

Així, el conjunt solució és:
\[
\boxed{
x\in \left[\frac{a-\sqrt{a^2+8}}{2},\;\frac{a+\sqrt{a^2+8}}{2}\right]
}
\]

\end{document}